% !TEX root = ../Main.tex
\chapter{一元二次不等式（不含参）}

\textbf{一元二次不等式}形如 $ax^2 + bx + c > 0$（或 $\ge, <, \le$），$a \ne 0$。本章目标：掌握"\textbf{看二次函数图像 $\to$ 读不等式解集}"这一核心方法——\textbf{代数问题转化为对抛物线开口、判别式、根位置的几何判断}。

\section{二次函数的三种表达式}

\state{二次函数的三种等价表达}{
    同一条抛物线可以用以下三种形式描述，每种突出不同的几何信息：
    \begin{itemize}
        \item \textbf{一般式} $y = ax^2 + bx + c$（$a \ne 0$）——突出系数 $a, b, c$；
        \item \textbf{顶点式} $y = a(x - h)^2 + k$——突出顶点 $(h, k)$；
        \item \textbf{两根式（因式式）} $y = a(x - x_1)(x - x_2)$——突出与 $x$ 轴的交点 $x_1, x_2$。
    \end{itemize}

    三种式子可以相互转换：
    \begin{itemize}
        \item \textbf{一般式 $\to$ 顶点式}：配方，$h = -\dfrac{b}{2a}$，$k = c - \dfrac{b^2}{4a}$；
        \item \textbf{一般式 $\to$ 两根式}：解方程 $ax^2 + bx + c = 0$，需判别式 $\Delta = b^2 - 4ac \ge 0$ 才有实根；
        \item \textbf{两根式 $\to$ 一般式}：直接展开。
    \end{itemize}
}

\state{判别式与根的情形}{
    对 $ax^2 + bx + c = 0$（$a \ne 0$），记 $\Delta = b^2 - 4ac$：
    \begin{itemize}
        \item $\Delta > 0$：\textbf{两不等实根} $x_{1,2} = \dfrac{-b \pm \sqrt{\Delta}}{2a}$，抛物线与 $x$ 轴交于两点；
        \item $\Delta = 0$：\textbf{两相等实根（重根）} $x = -\dfrac{b}{2a}$，抛物线与 $x$ 轴相切；
        \item $\Delta < 0$：\textbf{无实根}，抛物线与 $x$ 轴不相交。
    \end{itemize}
}

\section{解一元二次不等式的图像法}

\textbf{三步诊断}：
\begin{enumerate}
    \item \textbf{看开口}：$a > 0$ 向上、$a < 0$ 向下；
    \item \textbf{算 $\Delta$}：决定抛物线与 $x$ 轴有几个交点；
    \item \textbf{读图象}：$ax^2 + bx + c > 0$ $\iff$ 图象在 $x$ 轴上方——\textbf{"取上"}；$< 0$ $\iff$ 图象在 $x$ 轴下方——\textbf{"取下"}。
\end{enumerate}

\ex[$\Delta < 0$ 的情形]{
    解不等式 $x^2 - x + 12 > 0$。
}

\sol{
    $a = 1 > 0$（开口向上），$\Delta = (-1)^2 - 4 \cdot 1 \cdot 12 = 1 - 48 = -47 < 0$——\textbf{无实根}。

    抛物线开口向上、且始终在 $x$ 轴上方——\textbf{$y$ 恒正}。故不等式 $y > 0$ 对一切 $x \in \RR$ 成立。

    \textbf{解集}：$\RR = (-\infty, +\infty)$。
}

\begin{center}
\begin{tikzpicture}[>=Stealth, scale=0.8]
    \draw[->] (-2.5,0) -- (3.5,0) node[right] {$x$};
    \draw[->] (0,-0.5) -- (0,4) node[above] {$y$};
    \draw[thick, blue, domain=-2:3, samples=50, smooth] plot (\x, {\x*\x - \x + 0.5 + 0.5});
    \node at (1.5, 3) [right] {$y = x^2 - x + 12$};
    \node at (0.5, -0.3) [below] {\footnotesize 恒在 $x$ 轴上方};
\end{tikzpicture}
\end{center}

\ex[两根式直接读解]{
    解不等式 $(x + 4)(x - 3) > 0$。
}

\sol{
    左边已是\textbf{两根式}——两根 $x_1 = -4, x_2 = 3$，$a = 1 > 0$（展开后 $x^2 + x - 12$，系数为正）。

    开口向上、两根为 $-4, 3$：图象在 $x$ 轴上方 $\iff x < -4$ 或 $x > 3$。

    \textbf{解集}：$(-\infty, -4) \cup (3, +\infty)$。

    \textbf{口诀}："\textbf{大于取两边、小于取中间}"（针对开口向上的情形）。
}

\ex[$\Delta = 0$ 的情形]{
    解不等式 $(x - 7)^2 > 0$。
}

\sol{
    $(x - 7)^2 \ge 0$ 恒成立，且 $= 0 \iff x = 7$。故 $(x - 7)^2 > 0 \iff x \ne 7$。

    \textbf{解集}：$\{x \mid x \ne 7\} = (-\infty, 7) \cup (7, +\infty)$。

    从抛物线视角：$y = (x-7)^2$ 开口向上、与 $x$ 轴\textbf{相切}于 $x = 7$——除切点外 $y > 0$。
}

\rmkb{
    \textbf{九种组合一览表}——按 $a$ 正负与 $\Delta$ 三档交叉：

    \begin{center}
    \renewcommand{\arraystretch}{1.3}
    \begin{tabular}{c|c|c|c}
    & $\Delta > 0$ (两根 $x_1 < x_2$) & $\Delta = 0$ (重根 $x_0$) & $\Delta < 0$ \\ \hline
    $a > 0$，$y > 0$ & $x < x_1$ 或 $x > x_2$ & $x \ne x_0$ & $\RR$ \\
    $a > 0$，$y < 0$ & $x_1 < x < x_2$ & $\varnothing$ & $\varnothing$ \\
    $a > 0$，$y \ge 0$ & $x \le x_1$ 或 $x \ge x_2$ & $\RR$ & $\RR$
    \end{tabular}
    \end{center}

    $a < 0$ 时把不等式两侧同乘 $-1$（\textbf{注意方向翻转}）转化为 $a > 0$ 再查表。
}
